\section{Architecture et déploiement en production}

Le modèle sécurisé entraîné en développement (section 3.4) doit être déployé dans un environnement de production garantissant la disponibilité, la sécurité des accès, et la traçabilité complète des opérations. Cette section présente l'architecture de production Docker intégrant deux zones de sécurité critiques : la Zone 4 (sécurisation des accès et détection d'anomalies) et la Zone 5 (gouvernance, audit et monitoring).

\subsection{Architecture Docker de production}

L'architecture de production repose sur une infrastructure containerisée orchestrée par Docker Compose (\texttt{docker-compose.prod.yml}), avec Nginx comme reverse proxy frontal et deux instances d'API FastAPI servant les modèles baseline et sécurisé.

\paragraph{Stack Docker de production}

L'environnement de production déploie 8 services dans un réseau bridge isolé :

\begin{itemize}
    \item \textbf{2 APIs FastAPI} : \texttt{baseline-api} (port 8001) et \texttt{secured-api} (port 8002) servant les modèles respectifs avec limites de ressources (2 CPU, 2-4GB RAM)
    \item \textbf{Nginx} : Reverse proxy HTTP/HTTPS (ports 80/443) avec load balancing et terminaison SSL
    \item \textbf{PostgreSQL 15} : Base de données pour utilisateurs, logs d'audit, et IPs bloquées par le WAF
    \item \textbf{Redis 7} : Cache pour sessions JWT, compteurs de rate limiting, et prédictions fréquentes
    \item \textbf{Prometheus} : Collecte de métriques temps réel (requêtes/sec, latence, accuracy)
    \item \textbf{Grafana} : Dashboards de monitoring avec alerting configuré
    \item \textbf{Web Frontend} : Interface HTML/JS servie par Nginx pour soumission d'images
\end{itemize}

\paragraph{Configuration Docker Compose}

\begin{lstlisting}[language=yaml, caption=docker-compose.prod.yml (extrait)]
services:
  # API Securisee
  secured-api:
    build:
      context: ..
      dockerfile: docker/Dockerfile.secured
    ports:
      - "8002:8000"
    environment:
      - MODEL_TYPE=secured
      - ENABLE_DEFENSES=true
      - POSTGRES_HOST=postgres
      - REDIS_HOST=redis
      - JWT_SECRET=${JWT_SECRET}
    volumes:
      - ../models/secured:/app/models/secured:ro
    restart: unless-stopped
    deploy:
      resources:
        limits:
          cpus: '2'
          memory: 4G

  # Nginx reverse proxy
  nginx:
    image: nginx:alpine
    ports:
      - "80:80"
      - "443:443"
    volumes:
      - ../configs/nginx/nginx.conf:/etc/nginx/nginx.conf:ro
      - ../configs/nginx/ssl:/etc/nginx/ssl:ro
    depends_on:
      - baseline-api
      - secured-api

  # Base de donnees PostgreSQL
  postgres:
    image: postgres:15-alpine
    environment:
      - POSTGRES_USER=${POSTGRES_USER}
      - POSTGRES_PASSWORD=${POSTGRES_PASSWORD}
    volumes:
      - postgres-data-prod:/var/lib/postgresql/data
    ports:
      - "127.0.0.1:5434:5432"
    healthcheck:
      test: ["CMD-SHELL", "pg_isready"]
      interval: 10s

  # Cache Redis
  redis:
    image: redis:7-alpine
    command: redis-server --appendonly yes
    volumes:
      - redis-data-prod:/data
    ports:
      - "127.0.0.1:6379:6379"
\end{lstlisting}

Les services PostgreSQL et Redis exposent leurs ports uniquement sur localhost (127.0.0.1) pour éviter l'accès externe. Les APIs sont configurées avec une politique de restart automatique (\texttt{unless-stopped}) et des health checks pour garantir la disponibilité.

\paragraph{Nginx reverse proxy et routage}

Nginx (\texttt{docker/nginx.conf}) agit comme point d'entrée unique HTTPS, routant les requêtes selon l'endpoint :

\begin{itemize}
    \item \texttt{/predict/baseline} → \texttt{baseline-api:8000}
    \item \texttt{/predict/secured} → \texttt{secured-api:8000}
    \item \texttt{/security/*} → \texttt{secured-api:8000} (auth, WAF, anomalies)
    \item \texttt{/audit/*} → \texttt{secured-api:8000} (logs, stats, export)
    \item \texttt{/} → \texttt{web:80} (interface HTML statique)
\end{itemize}

Le proxy préserve l'IP réelle du client via les headers \texttt{X-Real-IP} et \texttt{X-Forwarded-For}, essentiels pour le rate limiting WAF par IP. Le certificat SSL assure le chiffrement TLS 1.3 avec suite de chiffrement forte.


\subsection{Zone 4 : Sécurisation en production}

La Zone 4 implémente quatre couches de défense protégeant l'API contre les accès non autorisés, les abus, et les comportements anormaux.

\paragraph{Web Application Firewall (WAF)}

Le WAF (\texttt{src/security/waf.py}) applique un rate limiting strict par IP et détecte les patterns d'attaque :

\begin{itemize}
    \item \textbf{Rate Limiting} : 100 requêtes/minute par IP (compteur Redis). Dépassement → HTTP 429, IP bloquée 15 minutes
    \item \textbf{Détection SQL Injection} : Regex détectant \texttt{UNION SELECT}, \texttt{DROP TABLE}, \texttt{--}. Match → HTTP 403, IP blacklistée
    \item \textbf{Détection XSS} : Regex détectant \texttt{<script>}, \texttt{onerror=}, \texttt{javascript:}. Match → requête rejetée
\end{itemize}

Les IPs bloquées sont stockées dans PostgreSQL avec timestamp, raison et durée du blocage.

\paragraph{Authentification JWT et contrôle d'accès RBAC}

L'authentification JWT (\texttt{src/security/auth.py}) utilise l'algorithme HS256. Le système RBAC définit trois rôles avec permissions distinctes :

\begin{lstlisting}[language=Python, caption=Définition des rôles et permissions]
class UserRole(Enum):
    ADMIN = "admin"      # Tous droits
    AGENT = "agent"      # Predictions + audit
    GUEST = "guest"      # Predictions uniquement

class Permission(Enum):
    PREDICT = "predict"
    VIEW_AUDIT = "view_audit"
    EXPORT_AUDIT = "export_audit"
    MANAGE_USERS = "manage_users"
    VIEW_METRICS = "view_metrics"

ROLE_PERMISSIONS = {
    UserRole.ADMIN: [
        Permission.PREDICT,
        Permission.VIEW_AUDIT,
        Permission.EXPORT_AUDIT,
        Permission.MANAGE_USERS,
        Permission.VIEW_METRICS
    ],
    UserRole.AGENT: [
        Permission.PREDICT,
        Permission.VIEW_AUDIT,
        Permission.VIEW_METRICS
    ],
    UserRole.GUEST: [
        Permission.PREDICT
    ]
}
\end{lstlisting}

Le flux d'authentification génère deux tokens : un access token (validité 15 minutes) contenant \texttt{user\_id}, \texttt{role} et \texttt{permissions}, et un refresh token (validité 7 jours) stocké dans Redis.

\paragraph{Détection d'anomalies comportementales}

Le détecteur d'anomalies (\texttt{src/security/anomaly\_detector.py}) analyse le comportement de chaque IP en temps réel sur plusieurs dimensions :

\begin{lstlisting}[language=Python, caption=Types d'anomalies et niveaux de risque]
class AnomalyType(Enum):
    SUSPICIOUS_PATTERN = "suspicious_pattern"
    BURST_ACTIVITY = "burst_activity"
    REPEATED_FAILURES = "repeated_failures"
    UNUSUAL_TIMING = "unusual_timing"
    CONFIDENCE_DROP = "confidence_drop"
    MODEL_SWITCHING = "model_switching"

class RiskLevel(Enum):
    LOW = "low"           # Surveiller
    MEDIUM = "medium"     # Alerte
    HIGH = "high"         # Bloquer
    CRITICAL = "critical" # Incident securite
\end{lstlisting}

Les features analysées incluent la fréquence de requêtes (fenêtre 1 minute), la diversité des endpoints, la taille moyenne des images, le ratio erreurs/succès, et le temps entre requêtes. Un score d'anomalie entre 0 et 1 est calculé, déclenchant des actions selon le seuil :

\begin{itemize}
    \item \textbf{Score 0.7-0.85} : Surveillance renforcée, log événement
    \item \textbf{Score 0.85-0.95} : Limitation temporaire (50 req/min)
    \item \textbf{Score >0.95} : Blocage automatique 1 heure, notification admin
\end{itemize}


\subsection{Zone 5 : Gouvernance et audit}

La Zone 5 garantit la traçabilité complète des opérations et l'observabilité temps réel du système.

\paragraph{Logs d'audit immuables}

L'audit logger (\texttt{src/monitoring/audit\_logger.py}) enregistre chaque événement critique dans PostgreSQL avec garantie d'immuabilité via hash chain SHA-256. Les types d'événements tracés sont :

\begin{lstlisting}[language=Python, caption=Types d'événements d'audit]
class EventType(Enum):
    PREDICTION = "prediction"
    ATTACK_DETECTED = "attack_detected"
    VALIDATION_FAILED = "validation_failed"
    API_ACCESS = "api_access"
    MODEL_LOADED = "model_loaded"
    SYSTEM_ERROR = "system_error"
    SECURITY_CHECK = "security_check"
    SECURITY_ALERT = "security_alert"
\end{lstlisting}

Chaque log contient :
\begin{itemize}
    \item \textbf{PREDICTION} : Timestamp, IP source (pseudonymisée), user\_id, model\_type, prediction, confidence, durée d'inférence, hash de l'image
    \item \textbf{ATTACK\_DETECTED} : IP, type d'attaque, payload, action prise
    \item \textbf{SECURITY\_ALERT} : Score d'anomalie, features, niveau de risque
\end{itemize}

Chaque log contient le hash SHA-256 du log précédent, créant une blockchain d'audit. Les IPs sont pseudonymisées avec HMAC-SHA256 (conformité RGPD).

\paragraph{Monitoring temps réel avec Prometheus et Grafana}

Prometheus (\texttt{configs/monitoring/prometheus.yml}) scrape les métriques des APIs toutes les 15 secondes. Les métriques exportées incluent :

\begin{itemize}
    \item \textbf{Compteurs} : \texttt{http\_requests\_total}, \texttt{ai\_model\_predictions\_total}, \texttt{waf\_blocks\_total}, \texttt{anomaly\_detections\_total}
    \item \textbf{Histogrammes} : \texttt{request\_duration\_seconds}, \texttt{model\_inference\_duration\_seconds} (P50, P95, P99)
    \item \textbf{Gauges} : \texttt{ai\_model\_accuracy}, \texttt{active\_jwt\_sessions}, \texttt{postgres\_connections}
\end{itemize}

Grafana affiche 4 dashboards pré-configurés :

\begin{enumerate}
    \item \textbf{Overview} : Requêtes/sec, latence P95, taux d'erreur, uptime
    \item \textbf{Security} : WAF blocks, anomalies, tentatives d'auth échouées, top IPs suspectes
    \item \textbf{Models} : Accuracy baseline vs secured, distribution prédictions, temps d'inférence
    \item \textbf{Audit} : Logs générés, top users, hash chain integrity
\end{enumerate}

Les règles d'alerte (\texttt{configs/monitoring/alert\_rules.yml}) se déclenchent sur des seuils critiques :

\begin{itemize}
    \item Taux d'erreur HTTP 5xx >10\% pendant 5 minutes → \texttt{HighErrorRate}
    \item Latence P95 >2 secondes pendant 10 minutes → \texttt{HighLatency}
    \item Accuracy <80\% sur fenêtre 1h → \texttt{ModelDegradation}
    \item Score d'anomalie moyen >0.8 pendant 15 minutes → \texttt{AnomalySpike}
\end{itemize}


\subsection{Déploiement et commandes}

Le déploiement en production utilise le Makefile pour simplifier les opérations courantes :

\begin{itemize}
    \item \texttt{make prod} : Build complet + déploiement (rebuild images)
    \item \texttt{make prod-fast} : Déploiement rapide sans rebuild
    \item \texttt{make health-prod} : Vérification health checks de tous les services
    \item \texttt{make logs-prod} : Logs en temps réel de tous les containers
    \item \texttt{make status-prod} : Statut des containers (running, exited, restart count)
    \item \texttt{make stop-prod} : Arrêt propre avec préservation des volumes
\end{itemize}

La commande \texttt{make prod} exécute séquentiellement : build des images Docker, démarrage des services avec \texttt{docker-compose up -d}, attente 30s pour PostgreSQL/Redis, vérification des health checks, et affichage du statut final.
